\documentclass[aspectratio=1610, 12pt]{beamer}

\usepackage[backend=biber,style=authoryear]{biblatex} % To RUN LOCALLY CHANGE BACKEND BIBER TO THE BUILT OPTIONS
\addbibresource{refs.bib} % Specify the path to your .bib file

\usepackage{graphicx}
\usepackage{amsmath}
\usepackage{media9}
\usepackage{multimedia} % for embedding videos
\usepackage{hyperref}
\usepackage{xcolor}
\usepackage{multimedia}
\usepackage{tikz}
\usetikzlibrary{calc}

% Theme configuration for UT Austin
\usetheme{default}

% Colors
\definecolor{White}{RGB}{255,255,255}

% Define the palette number (change this number to select a different palette)
\def\palette{1} % Change this number to select different palettes (1 to 8)

% Conditional setup for color palettes
\ifnum\palette=1
    % Palette 1: Navy Blue, Gray, and Dark Red
    \definecolor{MainColorTheme}{RGB}{46, 59, 78} % Navy Blue
    \definecolor{SecondaryColorTheme}{RGB}{125, 125, 125} % Gray
    \definecolor{TertiaryColorTheme}{RGB}{168, 50, 50} % Dark Red
\fi

\ifnum\palette=2
    % Palette 2: Dark Green, Gold, and Charcoal
    \definecolor{MainColorTheme}{RGB}{47, 79, 79} % Dark Green
    \definecolor{SecondaryColorTheme}{RGB}{218, 165, 32} % Gold
    \definecolor{TertiaryColorTheme}{RGB}{51, 51, 51} % Charcoal
\fi

\ifnum\palette=3
    % Palette 3: Midnight Blue, Silver, and Deep Teal
    \definecolor{MainColorTheme}{RGB}{27, 58, 87} % Midnight Blue
    \definecolor{SecondaryColorTheme}{RGB}{192, 192, 192} % Silver
    \definecolor{TertiaryColorTheme}{RGB}{0, 123, 127} % Deep Teal
\fi

\ifnum\palette=4
    % Palette 4: Forest Green, Slate Gray, and Burnt Orange
    \definecolor{MainColorTheme}{RGB}{59, 107, 53} % Forest Green
    \definecolor{SecondaryColorTheme}{RGB}{112, 128, 144} % Slate Gray
    \definecolor{TertiaryColorTheme}{RGB}{191, 87, 0} % Burnt Orange
\fi

\ifnum\palette=5
    % Palette 5: Royal Blue, Light Gray, and Burgundy
    \definecolor{MainColorTheme}{RGB}{65, 105, 225} % Royal Blue
    \definecolor{SecondaryColorTheme}{RGB}{211, 211, 211} % Light Gray
    \definecolor{TertiaryColorTheme}{RGB}{128, 0, 32} % Burgundy
\fi

\ifnum\palette=6
    % Palette 6: Dark Cyan, Taupe, and Burgundy
    \definecolor{MainColorTheme}{RGB}{0, 139, 139} % Dark Cyan
    \definecolor{SecondaryColorTheme}{RGB}{185, 165, 141} % Taupe
    \definecolor{TertiaryColorTheme}{RGB}{128, 0, 0} % Burgundy
\fi

\ifnum\palette=7
    % Palette 7: Oxford Blue, Light Steel Blue, and Amber
    \definecolor{MainColorTheme}{RGB}{0, 33, 71} % Oxford Blue
    \definecolor{SecondaryColorTheme}{RGB}{176, 196, 222} % Light Steel Blue
    \definecolor{TertiaryColorTheme}{RGB}{255, 191, 0} % Amber
\fi

\ifnum\palette=8
    % Palette 8: Steel Blue, Cool Gray, and Brick Red
    \definecolor{MainColorTheme}{RGB}{70, 130, 180} % Steel Blue
    \definecolor{SecondaryColorTheme}{RGB}{138, 141, 143} % Cool Gray
    \definecolor{TertiaryColorTheme}{RGB}{178, 34, 34} % Brick Red
\fi


% Set default color for titles and bullet points to Dark Orange (including the contents)
\setbeamercolor{frametitle}{fg=MainColorTheme}
\setbeamercolor{section in toc}{fg=MainColorTheme}

% Set bullet point symbols to Dark Orange, and text to Black
\setbeamercolor{item}{fg=MainColorTheme} % Bullet point symbols
\setbeamercolor{itemize/enumerate body}{fg=black} % Bullet point text
\setbeamercolor{itemize/enumerate subbody}{fg=black}

% Adjust the position of slide titles
\setbeamertemplate{frametitle}{
    \vspace{-0.7cm} % Adjust this value to move the title up or down
    \insertframetitle
}
% Customize figure caption format removing Figure
%\setbeamertemplate{caption}{\raggedright\insertcaption\par}

% Customize figure caption format with Dark Orange "Figure: ", numbering and caption size
\setbeamertemplate{caption}{
    \raggedright
    %\small
    \footnotesize 
    %\scriptsize 
    %\tiny
    \textcolor{MainColorTheme}{\insertcaptionname~\insertcaptionnumber:}~\insertcaption
}

% Customize reference slide caption format
\setbeamercolor{bibliography item}{fg=black}
\setbeamercolor{bibliography entry author}{fg=MainColorTheme}
\setbeamercolor{bibliography entry title}{fg=black}
\setbeamercolor{bibliography entry location}{fg=MainColorTheme}
\setbeamercolor{bibliography entry note}{fg=MainColorTheme}




% Title slide setup
%\setbeamertemplate{title page}{
%    % [T] aligns the tops of the columns, preventing the "push down" effect
%    \begin{columns}[T,onlytextwidth]
%            
%            \begin{column}{0.70\textwidth} % Left column
%               \vspace{.2cm}
%                \textcolor{MainColorTheme}{\textbf{\small MONTH 20XX}} \\
%                \vspace{1.5cm}
%                \textcolor{MainColorTheme}{\textbf{\Huge INSERT YOUR HEADLINE HERE UP TO 3 LINES}} \\
%                \vspace{.1cm}
%                \textcolor{MainColorTheme}{\rule{\linewidth}{0.1mm}} \\
%                \vspace{.2cm}
%                \textcolor{MainColorTheme}{Insert your subtitle or any additional description text here up to two lines of text or you can delete this text box} \\
%                \vspace{0.7cm}
%                \textcolor{MainColorTheme}{\textbf{\small Presenter or Speaker Name} \\ \small Position/Role, The University of Texas at Austin}
%            \end{column}
%
%            \begin{column}{0.3\textwidth} % Right column
%                % Use negative vspace to pull logos UP to the top edge
%                % Use positive vspace to push them DOWN
%                \vspace{0.1cm} 
%               \hfill
%                % The zero-width box allows the logos to expand leftwards freely
%                \makebox[0pt][r]{%
%                    % First Logo (NTU)
%                    %\includegraphics[width=5cm, keepaspectratio]{brand/ntu_cee_logo_navy.pdf}%
%                    %\includegraphics[width=4cm, keepaspectratio]{brand/DISC_NTU_full_names.png}%
%                    \includegraphics[width=5cm, keepaspectratio]{brand/DISC_NTU_full_names_vertical.png}%
%                    %\includegraphics[width=5cm, keepaspectratio]{brand/DISC_NTU_full_names_vertical_2.png}%
%                    %\includegraphics[width=5cm, keepaspectratio]{brand/DISC_NTU_full_names_vertical_3.png}%
%                    %\includegraphics[width=9cm, keepaspectratio]{brand/DISC_NTU_full_names_horizontal.png}%
%                    %\hspace{0.1cm} 
%                    % Second Logo (DISC)
%                    %\includegraphics[width=.93cm, keepaspectratio]{brand/logo_disc_0001_500px.png}%
%                    %\includegraphics[width=.93cm, keepaspectratio]{brand/DISC_fullname.png}%
%                }
%            \end{column}
%        
%    \end{columns}
%}

% Requires in preamble:
% \usepackage{tikz}
% \usetikzlibrary{calc}

% ---- Title slide setup (original layout, with background + new logo placement) ----
\setbeamertemplate{title page}{
\begin{tikzpicture}[remember picture,overlay]
    % Background image (full bleed)
    \node[anchor=center] at (current page.center) {%
        \includegraphics[width=\paperwidth,height=\paperheight]{brand/background_titlepage.png}%
    };
\end{tikzpicture}

% Foreground content (keep your original columns to preserve exact spacing)
\begin{columns}[T,onlytextwidth]

    % -------- Left column (TEXT) --------
    \begin{column}{0.70\textwidth}
        % Keep the original spacing, BUT remove the month here (we'll place it bottom-right)
        \vspace{.2cm}

        \vspace{1.5cm}
        \textcolor{MainColorTheme}{\textbf{\Huge INSERT YOUR HEADLINE HERE UP TO 3 LINES}}\\
        \vspace{.1cm}
        \textcolor{MainColorTheme}{\rule{\linewidth}{0.1mm}}\\
        \vspace{.2cm}
        \textcolor{MainColorTheme}{Insert your subtitle or any additional description text here up to two lines of text or you can delete this text box}\\
        \vspace{0.7cm}
        \textcolor{MainColorTheme}{\textbf{\small Presenter or Speaker Name}\\ \small Position/Role, Nanyang Technological University, Singapore}
    \end{column}

    % -------- Right column (empty; we place logos with overlay to control exact position) --------
    \begin{column}{0.30\textwidth}
        \vspace{0.1cm}
        % Intentionally left empty to avoid fighting with overlay positioning
    \end{column}

\end{columns}

% ---- Overlay: logos + month bottom-right (new design placement) ----
\begin{tikzpicture}[remember picture,overlay]

    % Top-left NTU logo (new design)
    \node[anchor=north west] at ([xshift=1.05cm,yshift=-0.65cm]current page.north west) {%
        \includegraphics[height=0.85cm]{brand/NTU_fullname.pdf}%
    };

    % Top-right DISC Lab logo (new design)
    \node[anchor=north east] at ([xshift=-1.05cm,yshift=-0.65cm]current page.north east) {%
        \includegraphics[height=0.85cm]{brand/DISC_fullname.pdf}%
    };

    % Month/Year at bottom-right
    \node[anchor=south east] at ([xshift=-1.05cm,yshift=0.75cm]current page.south east) {%
        \textcolor{MainColorTheme}{\textbf{\small MONTH 20XX}}%
    };

\end{tikzpicture}
}


%\setbeamertemplate{title page}{
%\begin{tikzpicture}[remember picture,overlay]
%    % --- Full-bleed background image ---
%    \node[anchor=center] at (current page.center) {%
%        \includegraphics[width=\paperwidth,height=\paperheight]{brand/background_titlepage.png}%
%    };
%    % --- Left vertical accent bar ---
%    \fill[MainColorTheme] (current page.north west) rectangle ([xshift=1.15cm]current page.south west);
%    % --- Thin accent line inside the bar ---
%    \fill[TertiaryColorTheme]
%       ([xshift=1.03cm,yshift=-0.8cm]current page.north west)
%        rectangle
%        ([xshift=1.06cm,yshift=0.8cm]current page.south west);
%    % --- Vertical date text (centered in the bar; no clipping) ---
%    \node[rotate=90, anchor=center, text=White] at ([xshift=0.58cm]current page.west) {%
%        \bfseries\small MONTH 20XX
%    };
%    % --- Logos band (top row, aligned + anchored) ---
%    \node[anchor=north west] at ([xshift=1.65cm,yshift=-0.55cm]current page.north west) {%
%        \includegraphics[height=0.85cm]{brand/NTU_fullname.png}%
%    };
%    \node[anchor=north east] at ([xshift=-1.20cm,yshift=-0.55cm]current page.north east) {%
%        \includegraphics[height=0.85cm]{brand/DISC_fullname.png}%
%    };
%    % --- Title block (reduced size + controlled line breaks) ---
%    \node[anchor=north west, text width=0.72\paperwidth] at
%        ([xshift=1.65cm,yshift=-2.05cm]current page.north west) {%
%        \raggedright
%        {\color{MainColorTheme}\bfseries\fontsize{38}{40}\selectfont
%            INSERT YOUR\\
%            HEADLINE\\
%            HERE
%        }\par\vspace{0.40cm}
%
%        {\color{MainColorTheme}\fontsize{15}{20}\selectfont
%            Insert your subtitle or any additional description text here up to two lines.
%        }\par
%    };
%    % --- Speaker block (separate node, pinned near bottom-left; never off-canvas) ---
%    \node[anchor=south west, text width=0.72\paperwidth] at
%        ([xshift=1.65cm,yshift=0.80cm]current page.south west) {%
%        \raggedright
%        {\color{MainColorTheme}\bfseries\fontsize{13}{16}\selectfont Presenter Name}\par
%        {\color{MainColorTheme}\fontsize{11.5}{14}\selectfont Position/Role, Nanyang Technological University}
%    };
%\end{tikzpicture}
%}



\newcommand{\setHeadlineText}[1]{
	\def\currentHeadlineText{#1}
}
% Header for subsequent slides with orange background
\setbeamertemplate{headline}{
    \begin{columns}[onlytextwidth]
        % Left Column (70% with Dark Orange Background)
        \begin{column}{0.7\textwidth}
            \vspace{-0.8cm}
            \begin{beamercolorbox}[wd=\paperwidth,ht=1.5cm,dp=0cm,center,sep=0pt]{palette primary}
                \usebeamercolor[bg]{palette primary}
                \vspace{.1cm}
                \hspace{-8.cm}
                %\includegraphics[height=.5cm]{brand/BGPR_white_horizontal.pdf} % Add the white logo here
                \includegraphics[height=.5cm]{brand/ntu_cee_logo_white.pdf} % Add the white logo here
                \hspace{0.1cm} \includegraphics[height=.55cm]{brand/DISC_fullname_white.pdf} % Add the white logo here
                %\hspace{0.1cm} \includegraphics[height=.5cm]{brand/logo_disc_0001_white_500px.png} % Add the white logo here
            \end{beamercolorbox}
        \end{column}
        % Right Column (30% with Light Orange Background)
        \begin{column}{0.3\textwidth}
            \vspace{-.8cm}
            \begin{beamercolorbox}[wd=\paperwidth,ht=1.5cm,dp=0cm,center,sep=0pt]{palette secondary}
                \vspace{.275cm}
                \hspace{-10.3cm}
                \textcolor{White}{\textbf{\tiny Digital Intelligence for Structures and Construction}}
                %\textcolor{White}{\textbf{\footnotesize \vspace{-0.1cm} \currentHeadlineText}}                
            \end{beamercolorbox}
        \end{column}
    \end{columns}
    \vspace{1.5cm}
}

% Define the beamer color boxes for custom colors
\setbeamercolor{palette primary}{bg=MainColorTheme,fg=White}
\setbeamercolor{palette secondary}{bg=SecondaryColorTheme,fg=White}

% General slide setup
\setbeamertemplate{navigation symbols}{}
\setbeamertemplate{footline}[frame number]




%HERE STARTS THE PRESENTATION SLIDES

\begin{document}

% Initial Slide
\begin{frame}[plain]    
    \titlepage
\end{frame}

% Table of Contents Slide
\begin{frame}{Table of Contents}
    \tableofcontents
\end{frame}

% Example Slide new section
\begin{frame}
    % Slide content goes here
    \begin{center}
        \textcolor{black}{\Huge Your Content Here}
    \end{center}
\end{frame}

% Example Slide new section
\begin{frame}
    % Slide content goes here
    \begin{center}
        \textcolor{TertiaryColorTheme}{\Huge Your Content Here}
    \end{center}
\end{frame}

\begin{frame}
    % Slide content goes here
    \begin{center}
        \textcolor{MainColorTheme}{\Huge Your Content Here}
    \end{center}
\end{frame}

% Section Divider Slide
\begin{frame}[plain]
    \vfill
    \centering
    \textcolor{MainColorTheme}{\Huge Section Title}
    \vfill
\end{frame}

% Slide with Bullet Points
\begin{frame}{Your Slide Title}
    \begin{itemize}
        \item Bullet point 1
        \item Bullet point 2
        \item Bullet point 3
        \item Bullet point 4
    \end{itemize}
\end{frame}

%  Slide with Two Columns (Bullet Points and an Image with Caption)
\begin{frame}{Your Slide Title}
    \begin{columns}
        % Left Column with Bullet Points
        \begin{column}{0.5\textwidth}
            \begin{itemize}
                \item Bullet point 1
                \item Bullet point 2
                \item Bullet point 3
            \end{itemize}
        \end{column}
        
        % Right Column with Image and Caption
        \begin{column}{0.5\textwidth}
            \begin{figure}
                \includegraphics[width=0.8\textwidth]{figures/image0.png}
                \caption{Your Image Caption}
            \end{figure}
        \end{column}
    \end{columns}
\end{frame}

% Bulletpoints at left and three images to the right
\begin{frame}{Slide Title}
    \begin{columns}
        % Left Column with Bullet Points
        \begin{column}{0.3\textwidth}
            \begin{itemize}
                \item Bullet point 1
                \item Bullet point 2
                \item Bullet point 3
                \item Bullet point 4
            \end{itemize}
        \end{column}
        
        % Right Column with Three Images Side by Side
        \begin{column}{0.7\textwidth}
            \begin{figure}
                \centering
                \includegraphics[width=0.32\textwidth]{figures/image1.png}
                \includegraphics[width=0.32\textwidth]{figures/image2.png}
                \includegraphics[width=0.32\textwidth]{figures/image3.png}
                \caption{Image captions can be placed below or omitted.}
            \end{figure}
        \end{column}
    \end{columns}
\end{frame}


% Image-Only Slide with Caption at the Bottom
\begin{frame}{Your Slide Title}
    \begin{figure}
        \includegraphics[height=0.45\textwidth]{figures/image1.png}
        \caption{Your Image Caption}
    \end{figure}
\end{frame}

% Full-Image Slide
\begin{frame}
    \begin{figure}
        \centering
        \includegraphics[width=\textwidth,height=0.9\textheight,keepaspectratio]{figures/sample_full_image.png}
    \end{figure}
\end{frame}


% Two Images Side by Side with Captions at the Bottom
\begin{frame}{Your Slide Title}
    \begin{columns}
        % Left Image
        \begin{column}{0.5\textwidth}
            \begin{figure}
                \includegraphics[width=.66\textwidth]{figures/image2.png}
                \caption{Caption for Image 1}
            \end{figure}
        \end{column}
        
        % Right Image
        \begin{column}{0.5\textwidth}
            \begin{figure}
                \includegraphics[width=.66\textwidth]{figures/image3.png}
                \caption{Caption for Image 2}
            \end{figure}
        \end{column}
    \end{columns}
\end{frame}

% Three Images Side by Side with Captions at the Bottom
\begin{frame}{Your Slide Title}
    \begin{columns}
        % Left Image
        \begin{column}{0.33\textwidth}
            \begin{figure}
                \includegraphics[width=\textwidth]{figures/image4.png}
                \caption{Caption for Image 1}
            \end{figure}
        \end{column}
        
        % Middle Image
        \begin{column}{0.33\textwidth}
            \begin{figure}
                \includegraphics[width=\textwidth]{figures/image5.png}
                \caption{Caption for Image 2}
            \end{figure}
        \end{column}
        
        % Right Image
        \begin{column}{0.33\textwidth}
            \begin{figure}
                \includegraphics[width=\textwidth]{figures/image6.png}
                \caption{Caption for Image 3}
            \end{figure}
        \end{column}
    \end{columns}
\end{frame}

% Three images as one figure
\begin{frame}{Slide Title}
    \begin{figure}
        \centering
        \includegraphics[width=0.32\textwidth]{figures/image1.png}
        \includegraphics[width=0.32\textwidth]{figures/image2.png}
        \includegraphics[width=0.32\textwidth]{figures/image3.png}
        \caption{Image captions can be placed below or omitted.}
    \end{figure}
\end{frame}

% Figure with 2 rows of 3 images
\begin{frame}{Slide Title}
    \begin{figure}
        \centering
        % First Row of Images
        \includegraphics[width=0.2\textwidth]{figures/image1.png} \hspace{0.1cm}
        \includegraphics[width=0.2\textwidth]{figures/image2.png} \hspace{0.1cm}
        \includegraphics[width=0.2\textwidth]{figures/image3.png} 
        \vspace{0.1cm} % Space between the rows        
        
        % Second Row of Images
        \includegraphics[width=0.2\textwidth]{figures/image4.png} \hspace{0.1cm}
        \includegraphics[width=0.2\textwidth]{figures/image5.png} \hspace{0.1cm}
        \includegraphics[width=0.2\textwidth]{figures/image6.png}
        \caption{Caption for the entire figure, describing all images}
    \end{figure}
\end{frame}

% Image at top, four images at bottom
\begin{frame}{Slide Title}
    \begin{figure}
        % Top Image
        \centering
        \includegraphics[height=0.25\textheight]{figures/image0.png}
        \caption{Top Image Caption (optional)}
    \end{figure}

    %\vspace{0.5cm} % Space between the top image and the bottom row of images

    \begin{figure}
        % Bottom Images Side by Side
        \centering
        \includegraphics[width=0.23\textwidth]{figures/image1.png}
        \includegraphics[width=0.23\textwidth]{figures/image2.png}
        \includegraphics[width=0.23\textwidth]{figures/image3.png}
        \includegraphics[width=0.23\textwidth]{figures/image4.png}
        \caption{Bottom Images Caption (optional)}
    \end{figure}
\end{frame}


% Video Slide
\begin{frame}{Video Demonstration}
    \begin{center}
        \includemedia[
            width=0.5\linewidth,
            height=0.3\linewidth,
            activate=onclick,
            addresource=videos/sample-video.mp4,
            flashvars={
               source=videos/sample-video.mp4
              &autoPlay=true
              &loop=true
            }
        ]{}{VPlayer.swf}
        \vspace{0.5cm}
        \textcolor{black}{A brief demonstration video.}
    \end{center}
\end{frame}

\begin{frame}{Video Demonstration}
    \frametitle{Video Demonstration}
    \href{run:videos/sample-video.mp4}{\beamergotobutton{Play Video}}
    \newline
    \textbf{A brief demonstration video.}
\end{frame}

\begin{frame}{Video Example}
    \movie[width=0.6\textwidth,height=0.4\textwidth,showcontrols]{}{videos/sample-video.mp4}
\end{frame}

% Example Table Slide
\begin{frame}{Table Slide Example}
    \begin{table}
        \centering
        \begin{tabular}{|c|c|c|}
            \hline
            \textbf{Column 1} & \textbf{Column 2} & \textbf{Column 3} \\
            \hline
            Row 1, Col 1 & Row 1, Col 2 & Row 1, Col 3 \\
            \hline
            Row 2, Col 1 & Row 2, Col 2 & Row 2, Col 3 \\
            \hline
            Row 3, Col 1 & Row 3, Col 2 & Row 3, Col 3 \\
            \hline
        \end{tabular}
        \caption{Example of a Table in a Slide}
    \end{table}
\end{frame}

% Image to a Table in LaTeX
\begin{frame}{Table with an Image}
    \begin{table}
        \centering
        \begin{tabular}{|c|c|c|}
            \hline
            \textbf{Column 1} & \textbf{Column 2} & \textbf{Column 3} \\
            \hline
            Row 1, Col 1 & \includegraphics[width=1.cm]{figures/image0.png} & Row 1, Col 3 \\
            \hline
            Row 2, Col 1 & Row 2, Col 2 & Row 2, Col 3 \\
            \hline
            Row 3, Col 1 & Row 3, Col 2 & Row 3, Col 3 \\
            \hline
        \end{tabular}
        \caption{Table with an Image in One Cell}
    \end{table}
\end{frame}

% Data/Statistics Slide
\begin{frame}{Key Statistics}
    \begin{itemize}
        \item \textbf{Stat 1:} 85\% of participants found the program beneficial.
        \item \textbf{Stat 2:} The average improvement was 25\% over six months.
        \item \textbf{Stat 3:} 75\% of users reported higher satisfaction.
    \end{itemize}
    \vspace{1cm}
    \centering
    \textcolor{MainColorTheme}{\large Data Source: \textit{Your Data Source}}
\end{frame}

% Chart/Graph Slide
\begin{frame}{Performance Over Time}
    \begin{figure}
        \centering
        \includegraphics[width=.7\textwidth]{figures/sample_chart.jpg}
        \caption{Performance trends across the quarters.}
    \end{figure}
\end{frame}

% Process/Timeline Slide
\begin{frame}{Project Timeline}
    \begin{center}
        \begin{tabular}{|c|c|c|c|}
            \hline
            \textbf{Phase 1} & \textbf{Phase 2} & \textbf{Phase 3} & \textbf{Phase 4} \\
            \hline
            January & March & May & July \\
            Planning & Development & Testing & Launch \\
            \hline
        \end{tabular}
    \end{center}
    \vspace{0.5cm}
    \centering
    \textcolor{MainColorTheme}{\large Key Milestones in the Project Lifecycle}
\end{frame}


% Adding Footnotes in Beamer Slides
\begin{frame}{Slide with Footnotes}
    Some important information is presented here.\footnote{This is a footnote explaining something important.}
    
    \begin{itemize}
        \item First point
        \item Second point\footnote{This footnote explains the second point.}
    \end{itemize}
\end{frame}

% Text and Quote Slide
\begin{frame}{Key Point}
    \begin{quote}
        \textcolor{MainColorTheme}{\large "This is a key quote that supports the presentation. It can be something profound, insightful, or important."}
    \end{quote}
    \hfill \\
    \raggedleft \textcolor{black}{\small -- Author Name or Source}
\end{frame}

% Quote or Testimonial Slide
\begin{frame}{Testimonial}
    \begin{quote}
        \textcolor{MainColorTheme}{\large "This service has greatly improved our efficiency and productivity. Highly recommend!"}
    \end{quote}
    \vspace{1cm}
    \raggedleft \textcolor{black}{\small -- Jane Doe, CEO of Example Corp.}
\end{frame}



% Cite Sources in Your Document
\begin{frame}{Citing Sources in LaTeX}
    This is a statement that requires citation.~\cite{author2023}
\end{frame}

% Thank You/Conclusion Slide
\begin{frame}%[plain]
    \vfill
    \centering
    \textcolor{MainColorTheme}{\Huge Thank you!}
    \vfill
    \textcolor{black}{\large Any questions?}
\end{frame}

% Q&A Slide
\begin{frame}[plain]
    \vfill
    \centering
    \textcolor{MainColorTheme}{\Huge Questions?}
    \vfill
\end{frame}


% Contact Information Slide
\begin{frame}{Contact Information}
    \textcolor{MainColorTheme}{\large Name:} \textcolor{black}{Name LastName} \\
    \textcolor{MainColorTheme}{\large Email:} \textcolor{black}{name.lastname@utexas.edu} \\
    \textcolor{MainColorTheme}{\large Phone:} \textcolor{black}{+123 456 7890} \\
    \textcolor{MainColorTheme}{\large Website:} \textcolor{black}{www.name-lastname.com} \\
    \vspace{1cm}
    \textcolor{MainColorTheme}{\large Social Media:}
    \textcolor{black}{
        \href{https://www.linkedin.com/in/username}{LinkedIn} \quad
        \href{https://twitter.com/username}{Twitter} \quad
        \href{https://github.com/username}{GitHub}
    }
\end{frame}


% Add a Bibliography Section
\begin{frame}{References}
    %\frametitle{References}
    \printbibliography
\end{frame}

\end{document}
